% Section 5.3, from lectures/L05/appendix/c_class_templates.md.

\section{Class Templates}
\label{c5:sec:class}\label{c5:app:c}
Class templates allow us to define generic data structures, where the type is determined at
compile time.

This enables reusable and efficient code without run-time overhead.

\subsection{Example: a generic container}
\label{c5:sec:container}
In this section, we implement a simple container similar to \code{std::vector}:

\begin{cppcode}
container::Vector<T>
\end{cppcode}

\noindent The purpose is to:
\begin{itemize}
  \item Understand how class templates work.
  \item Understand how memory management affects embedded systems.
  \item Use compile-time constraints for safety.
  \item See compile-time polymorphism in practice.
\end{itemize}

\subsection{What is a class template?}
\label{c5:sec:whatisclasstemplate}
A class template is a blueprint for a class where the type is not fixed until compilation.

Example:

\begin{cppcode}
namespace container
{
template<typename T>
class Vector
{
public:
    Vector() noexcept
        : myData{nullptr}
        , mySize{}
    {}
private:
    T* myData;
    std::size_t mySize;
};
} // namespace container
\end{cppcode}

\noindent When used:

\begin{cppcode}
container::Vector<int> a{};
container::Vector<double> b{};
\end{cppcode}

\noindent The compiler generates two separate types:

\begin{cppcode}
namespace container
{
// Vector specialization for int.
class Vector<int>
{
public:
    Vector() noexcept
        : myData{nullptr}
        , mySize{}
    {}
private:
    int* myData;
    std::size_t mySize;
};

// Vector specialization for double.
class Vector<double>
{
public:
    Vector() noexcept
        : myData{nullptr}
        , mySize{}
    {}
private:
    double* myData;
    std::size_t mySize;
};
} // namespace container
\end{cppcode}

\noindent This is compile-time polymorphism.

\subsection{Example implementation}
\label{c5:sec:vectorimpl}
Below is a simplified implementation of a vector-like container.

Key features:
\begin{itemize}
  \item Dynamic resizing.
  \item Push-back functionality.
  \item Range-based iteration support.
\end{itemize}

\begin{cppcode}
#pragma once

#include <cstddef>
#include <cstdlib>
#include <type_traits>

namespace container
{
template<typename T>
class Vector
{
    // Generate a compiler error if the given type isn't trivially copyable or destructible,
    // since std::realloc() and std::free() are used in this implementation.
    static_assert(std::is_trivially_copyable<T>::value && std::is_trivially_destructible<T>::value,
        "Vector type T must be of trivial type!");
public:
    Vector() noexcept
        : myData{nullptr}
        , mySize{}
    {
        // Initialized empty vector.
    }

    template<typename... Args>
    Vector(const Args& ...args) noexcept
        : Vector()
    {
        // Try to resize the vector to fit the given number of arguments.
        // Copy the arguments to the vector on success.
        if (resize(sizeof...(args)))
        {
            std::size_t i{};
            ((myData[i++] = args), ...);
        }
    }

    ~Vector() noexcept
    {
        // Release allocated resources before deletion.
        clear();
    }

    T& operator[](const std::size_t index) noexcept
    {
        // Return a reference to the element at given index.
        return myData[index];
    }
    const T& operator[](const std::size_t index) const noexcept
    {
        // Return a read-only reference to the element at given index.
        return myData[index];
    }

    T* begin() noexcept
    {
        // Return a pointer to the beginning of the data field.
        return myData;
    }
    const T* begin() const noexcept
    {
        // Return a read-only pointer to the beginning of the data field.
        return myData;
    }

    T* end() noexcept { return nullptr != myData ? myData + mySize : nullptr; }
    const T* end() const noexcept { return nullptr != myData ? myData + mySize : nullptr; }

    [[nodiscard]] const T* data() const noexcept
    {
        // Return a read-only pointer to the data field.
        return myData;
    }

    [[nodiscard]] std::size_t size() const noexcept
    {
        // Return the size of the vector.
        return mySize;
    }

    [[nodiscard]] bool empty() const noexcept
    {
        // Return true if the vector is empty.
        return 0U == mySize;
    }

    void clear() noexcept
    {
        // Free allocated resources, then reset the vector parameters.
        std::free(myData);
        myData = nullptr;
        mySize = 0U;
    }

    [[nodiscard]] bool resize(const std::size_t size) noexcept
    {
        // Clear the vector if the requested size is 0.
        if (0U == size)
        {
            clear();
            return true;
        }

        // Try to reallocate the data field to the requested size, return false on failure.
        T* data{static_cast<T*>(std::realloc(myData, sizeof(T) * size))};
        if (nullptr == data) { return false; }
        myData = data;

        // Initialize new elements (if any).
        // Update the vector size, then return true to indicate success.
        for (std::size_t i{mySize}; i < size; ++i) { myData[i] = {}; }
        mySize = size;
        return true;
    }

    [[nodiscard]] bool pushBack(const T& element) noexcept
    {
        const std::size_t newSize{mySize + 1U};
        // Try to reallocate the data field to fit one more element, return false on failure.
        T* data{static_cast<T*>(std::realloc(myData, sizeof(T) * newSize))};
        if (nullptr == data) { return false; }

        // Push the new element to the back of the vector, then return true to indicate success.
        myData = data;
        myData[mySize++] = element;
        return true;
    }

    Vector(const Vector&)            = delete; // No copy constructor.
    Vector(Vector&&)                 = delete; // No move constructor.
    Vector& operator=(const Vector&) = delete; // No copy assignment.
    Vector& operator=(Vector&&)      = delete; // No move assignment.

private:
    /** Pointer to dynamically allocated data field. */
    T* myData;

    /** Size of the data field in number of elements. */
    std::size_t mySize;
};
} // namespace container
\end{cppcode}

\noindent The above implementation means that the class template \code{container::Vector} behaves
in many ways like \code{std::vector}, including initialization with an argument list:

\begin{cppcode}
container::Vector<int> vec1{1, 3, 5, 7, 9};
container::Vector<double> vec2{0.5, 1.5, 2.5, 3.5, 4.5};
\end{cppcode}

\subsection{Placement of the implementation in header files}
\label{c5:sec:headers}
Since \code{container::Vector} is a template, this affects how the implementation must be placed
in the source code.

The compiler must see the full definition of a template when it instantiates it, for example when
we write:

\begin{cppcode}
container::Vector<int>
\end{cppcode}

\noindent This means that both the class declaration and the method implementations must be
visible where the template is used.

There are two common solutions.

\subsubsection{a) Everything in the same header file}
The most common approach is to place both declaration and implementation in the same header file,
for example:

\begin{console}
include/container/vector.hpp
\end{console}

\noindent Contains:

\begin{cppcode}
template<typename T>
class Vector { ... };

template<typename T>
bool Vector<T>::pushBack(const T& element)
{
    ...
}
\end{cppcode}

\noindent This works because the compiler sees all code when the template is instantiated.

\subsubsection{b) Header + implementation header}
In larger projects, it is common to separate declaration and implementation into two files, but
include the implementation at the end of the header file.

Structure:

\begin{console}
include/container/vector.hpp
include/container/impl/vector_impl.hpp
\end{console}

\noindent And at the end of \file{vector.hpp}:

\begin{cppcode}
#include "container/impl/vector_impl.hpp"
\end{cppcode}

\noindent This gives the same effect as option a), but provides better structure and readability.

\subsubsection{Why doesn't a \code{.cpp} file work?}
If the implementation is placed in a regular \code{.cpp} file, the compiler cannot see the code
when the template is instantiated in other files. This often leads to linker errors such as:

\begin{console}
undefined reference to ...
\end{console}

\noindent Templates must therefore be implemented in header files (directly or indirectly via
include).

\subsection{Embedded perspective}
\label{c5:sec:embedded}
In embedded systems, we need to consider the following:

\subsubsection{Heap usage}
\begin{itemize}
  \item Dynamic allocation may be forbidden.
  \item Fragmentation can occur.
\end{itemize}

\subsubsection{Binary size}
\begin{itemize}
  \item Each instantiation of \code{Vector<T>} generates new code.
  \item \code{Vector<int>} $\neq$ \code{Vector<double>}.
\end{itemize}

\subsubsection{Exception handling}
\begin{itemize}
  \item In many embedded systems, exceptions are disabled.
  \item Therefore, \code{noexcept} is used.
\end{itemize}

\subsection{Summary}
\label{c5:sec:classsummary}
In this section, we have seen:
\begin{itemize}
  \item How to create a class template.
  \item How to handle manual memory allocation.
  \item How copy/move semantics can be handled or restricted.
  \item Why template implementations must reside in header files.
  \item How templates affect embedded systems.
\end{itemize}
\code{container::Vector<T>} is a clear example of:
\begin{itemize}
  \item Compile-time polymorphism.
  \item Generic programming.
  \item Performance-optimized C++.
  \item Full control over memory and lifecycle.
\end{itemize}
